\cleardoublepage
\chapter{FIXTURE EVALUASI PAKAR}

Lampiran ini menyajikan 10 \textit{fixture} evaluasi yang digunakan dalam validasi pakar (\textit{expert evaluation}) sebagaimana diuraikan pada Subbab~\ref{sec:validasi-pakar-jawaban}. Setiap \textit{fixture} mewakili satu pertanyaan beserta \textit{ground truth} yang mencakup \textit{relevant nodes} dan \textit{expected path} di \textit{knowledge graph}. \textit{Field} \texttt{answer} dan \texttt{context} dikecualikan dari tampilan ini untuk mempertahankan objektivitas penilaian. Tabel~\ref{tbl:fixture-summary} merangkum distribusi 10 \textit{fixture} tersebut.

\begin{table}[h]
\centering
\caption{Ringkasan 10 Fixture Evaluasi Pakar}
\label{tbl:fixture-summary}
{\small
\setlength{\tabcolsep}{4pt}
\begin{tabular}{|c|p{3.3cm}|p{2cm}|p{2.6cm}|p{1.6cm}|c|}
\hline
\textbf{No} & \textbf{ID} & \textbf{Tipe} & \textbf{Resource Utama} & \textbf{Scope} & \textbf{M-H} \\
\hline
1 & \textit{deployment\_basic}            & konseptual    & Deployment              & Namespaced & Ya \\
2 & \textit{service\_types}               & konseptual    & Service                 & Namespaced & Ya \\
3 & \textit{scale\_existing \_deployment}  & lanjutan      & Deployment              & Namespaced & Tidak \\
4 & \textit{ingress\_service \_pod}        & relasional    & Ingress                 & Namespaced & Ya \\
5 & \textit{hpa\_deployment \_pod}         & relasional    & HorizontalPodAutoscaler & Namespaced & Ya \\
6 & \textit{cronjob\_backup}              & generasi YAML & CronJob                 & Namespaced & Ya \\
7 & \textit{networkpolicy \_deny\_all}     & generasi YAML & NetworkPolicy           & Namespaced & Tidak \\
8 & \textit{statefulset\_with \_pvc}       & generasi YAML & StatefulSet             & Namespaced & Ya \\
9 & \textit{deployment\_vs\_statefulset \_comparison} & skenario nyata & Deployment      & Namespaced    & Ya \\
10 & \textit{kubectl\_find \_pods\_with \_env} & perintah     & Pod                    & Namespaced & Tidak \\
\hline
\end{tabular}
}
\end{table}

% ─────────────────────────────────────────────────────────────────────────────
\section{Fixture B.1 — \texttt{deployment\_basic} (Konseptual)}
% ─────────────────────────────────────────────────────────────────────────────

\textit{Fixture} ini menguji pemahaman sistem terhadap mekanisme \textit{RollingUpdate} Deployment: bagaimana Pod diperbarui saat \textit{image} berubah dan apakah terjadi \textit{downtime}.

\lstinputlisting[
    language=json,
    frame=single,
    captionpos=t,
    caption={\texttt{fixture-deployment-basic.json} — Konseptual: Pembaruan Image Deployment},
    label={lst:fixture-deployment-basic}
]{listings/fixture-deployment-basic.json}

% ─────────────────────────────────────────────────────────────────────────────
\section{Fixture B.2 — \texttt{service\_types} (Konseptual)}
% ─────────────────────────────────────────────────────────────────────────────

\textit{Fixture} ini menguji kemampuan sistem membandingkan dua mekanisme \textit{expose} aplikasi: Service tipe LoadBalancer versus Ingress, beserta pertimbangan \textit{trade-off}-nya.

\lstinputlisting[
    language=json,
    frame=single,
    captionpos=t,
    caption={\texttt{fixture-service-types.json} — Konseptual: Service LoadBalancer vs.\ Ingress},
    label={lst:fixture-service-types}
]{listings/fixture-service-types.json}

% ─────────────────────────────────────────────────────────────────────────────
\section{Fixture B.3 — \texttt{scale\_existing\_deployment} (Lanjutan)}
% ─────────────────────────────────────────────────────────────────────────────

\textit{Fixture} bertipe \textit{followup} ini mengukur kemampuan sistem memanfaatkan konteks percakapan sebelumnya (pembuatan Deployment nginx) untuk menjawab pertanyaan lanjutan tentang perubahan jumlah replika.

\lstinputlisting[
    language=json,
    frame=single,
    captionpos=t,
    caption={\texttt{fixture-scale-existing-deployment.json} — Lanjutan: Skalakan Replika Deployment},
    label={lst:fixture-scale-deployment}
]{listings/fixture-scale-existing-deployment.json}

% ─────────────────────────────────────────────────────────────────────────────
\section{Fixture B.4 — \texttt{ingress\_service\_pod} (Relasional)}
% ─────────────────────────────────────────────────────────────────────────────

\textit{Fixture} ini menguji penelusuran relasi multi-hop: alur traffic dari Ingress melalui Service hingga ke Pod, dan identifikasi titik kegagalan pada error 502.

\lstinputlisting[
    language=json,
    frame=single,
    captionpos=t,
    caption={\texttt{fixture-ingress-service-pod.json} — Relasional: Alur Traffic Ingress→Service→Pod},
    label={lst:fixture-ingress-service}
]{listings/fixture-ingress-service-pod.json}

% ─────────────────────────────────────────────────────────────────────────────
\section{Fixture B.5 — \texttt{hpa\_deployment\_pod} (Relasional)}
% ─────────────────────────────────────────────────────────────────────────────

\textit{Fixture} ini menguji kemampuan sistem menjelaskan siklus kerja HPA: alur dari pembacaan metrik CPU hingga pembaruan \texttt{spec.replicas} di Deployment, dan diagnosis ketika HPA tidak berfungsi.

\lstinputlisting[
    language=json,
    frame=single,
    captionpos=t,
    caption={\texttt{fixture-hpa-deployment-pod.json} — Relasional: HPA dan Penskalaan Otomatis},
    label={lst:fixture-hpa-deployment}
]{listings/fixture-hpa-deployment-pod.json}

% ─────────────────────────────────────────────────────────────────────────────
\section{Fixture B.6 — \texttt{cronjob\_backup} (Generasi YAML)}
% ─────────────────────────────────────────────────────────────────────────────

\textit{Fixture} generasi YAML ini mengukur kemampuan sistem menghasilkan manifes CronJob yang valid dengan jadwal \textit{cron} yang tepat dan konfigurasi \textit{job template} yang benar.

\lstinputlisting[
    language=json,
    frame=single,
    captionpos=t,
    caption={\texttt{fixture-cronjob-backup.json} — Generasi YAML: CronJob Backup Database},
    label={lst:fixture-cronjob}
]{listings/fixture-cronjob-backup.json}

% ─────────────────────────────────────────────────────────────────────────────
\section{Fixture B.7 — \texttt{networkpolicy\_deny\_all} (Generasi YAML)}
% ─────────────────────────────────────────────────────────────────────────────

\textit{Fixture} ini menguji generasi NetworkPolicy dengan \textit{empty podSelector} untuk memblokir semua \textit{ingress traffic} ke namespace—pola keamanan yang umum tetapi memerlukan pemahaman semantik \texttt{podSelector: \{\}}.

\lstinputlisting[
    language=json,
    frame=single,
    captionpos=t,
    caption={\texttt{fixture-networkpolicy-deny-all.json} — Generasi YAML: NetworkPolicy Deny-All},
    label={lst:fixture-networkpolicy}
]{listings/fixture-networkpolicy-deny-all.json}

% ─────────────────────────────────────────────────────────────────────────────
\section{Fixture B.8 — \texttt{statefulset\_with\_pvc} (Generasi YAML)}
% ─────────────────────────────────────────────────────────────────────────────

\textit{Fixture} ini menguji generasi StatefulSet dengan \texttt{volumeClaimTemplates}—field yang spesifik untuk StatefulSet dan tidak ada di Deployment—untuk menguji kemampuan sistem membedakan kedua \textit{workload} tersebut.

\lstinputlisting[
    language=json,
    frame=single,
    captionpos=t,
    caption={\texttt{fixture-statefulset-with-pvc.json} — Generasi YAML: StatefulSet dengan PVC},
    label={lst:fixture-statefulset}
]{listings/fixture-statefulset-with-pvc.json}

% ─────────────────────────────────────────────────────────────────────────────
\section{Fixture B.9 — \texttt{deployment\_vs\_statefulset\_comparison} (Skenario Nyata)}
% ─────────────────────────────────────────────────────────────────────────────

\textit{Fixture} bertipe \textit{realworld} ini berasal dari pertanyaan Stack Overflow. Pertanyaan tentang perbedaan skema antara Deployment dan StatefulSet menguji kemampuan sistem melakukan perbandingan lintas \textit{resource}, sebuah pola pertanyaan konseptual dengan dua \textit{root node} sekaligus.

\lstinputlisting[
    language=json,
    frame=single,
    captionpos=t,
    caption={\texttt{fixture-deployment-vs-statefulset.json} — Skenario Nyata: Perbandingan Deployment vs.\ StatefulSet},
    label={lst:fixture-deployment-vs-statefulset}
]{listings/fixture-deployment-vs-statefulset.json}

% ─────────────────────────────────────────────────────────────────────────────
\section{Fixture B.10 — \texttt{kubectl\_find\_pods\_with\_env} (Perintah)}
% ─────────────────────────────────────────────────────────────────────────────

\textit{Fixture} bertipe \textit{command} ini menguji kemampuan sistem memberikan perintah \texttt{kubectl} yang tepat untuk mencari \textit{environment variable} di seluruh Pod—pertanyaan operasional yang membutuhkan pemahaman struktur \texttt{Container.env} dan \texttt{envFrom}.

\lstinputlisting[
    language=json,
    frame=single,
    captionpos=t,
    caption={\texttt{fixture-kubectl-find-pods-with-env.json} — Perintah: Pencarian Environment Variable},
    label={lst:fixture-kubectl-env}
]{listings/fixture-kubectl-find-pods-with-env.json}
